\documentclass[../main.tex]{subfiles}

\begin{document}
\begin{CJK*}{UTF8}{gbsn}

\section*{Exercise 5}
Fix $n \geqslant 1$, $\tau \in (0,1)$.
For any vector $v \in \mathbb{R}^n$,
define $\text{Quantile}(v; \tau) = v_{(\lceil \tau n \rceil)}$,
where $v_{(1)} \leqslant v_{(2)} \leqslant \ldots \leqslant v_{(n)}$ are the order statistics of $v$.
Under what condition is it true that, for all $v \in \mathbb{R}^n$,
$\text{Quantile}(v; 1- \tau) = -\text{Quantile}(-v; \tau)$?

\paragraph{Solution}
The order statistics of $-v$ are given by $(-v)_{(i)} = -v_{(n-i+1)}$ for $i = 1, \ldots, n$.
Therefore, $\text{Quantile}(v; 1- \tau) = -\text{Quantile}(-v; \tau)$
holds
for a vector $v \in \mathbb{R}^n$
if and only if 

\begin{equation*}
  v_{(\lceil (1-\tau) n \rceil)} = -((-v)_{(\lceil \tau n \rceil)})
  = -(-v_{(n-\lceil \tau n \rceil + 1)})
  = v_{(n-\lceil \tau n \rceil + 1)}.
\end{equation*}
To make this work for all $v \in \mathbb{R}^n$, the condition we must pose is that
\begin{equation*}
\lceil (1-\tau) n \rceil = n-\lceil \tau n \rceil + 1.
\end{equation*}
This can happen for instance when $\tau = 1/2$ and $n$ is odd.
////
\end{CJK*}
\end{document}